\documentclass[10pt,a4paper]{amsart}
\usepackage[margin=19mm]{geometry}
\usepackage{amsmath,amssymb,mathtools}
\usepackage[scr=rsfs,cal=euler]{mathalfa}
\usepackage{xfrac}
\usepackage[unicode,hidelinks,bookmarks=false]{hyperref}
\newcommand{\ie}{i.e.\ }
\newcommand{\cf}{cf.\ }
\newcommand{\resp}{respectively\ }
\newcommand{\Subsection}{\subsection}
\providecommand{\nobreakdash}{\nobreak\hbox{-}}
\setlength{\emergencystretch}{2em}
\multlinegap0pt
\usepackage{amssymb}
\usepackage{tikz}
\newtheorem{proposition}{Proposition}
\newtheorem{theorem}{Theorem}
\newtheorem{lemma}{Lemma}
\usepackage{cleveref}
\crefname{proposition}{Proposition}{Propositions}
\crefname{theorem}{Theorem}{Theorems}
\crefname{lemma}{Lemma}{Lemmas}
\DeclareMathOperator{\proj}{\Pi}
\title{Bilateral facial reduction: qualification-free subdifferential calculus and exact duality}
\author{Matthew S. Scott}
\date{Source: arXiv:2510.12244v5 (CC BY 4.0)}
\begin{document}
\maketitle

\noindent\emph{Source and licence.} Bilateral facial reduction: qualification-free subdifferential calculus and exact duality. Version arXiv:2510.12244v5 (CC BY 4.0), distributed by arXiv under the Creative Commons Attribution 4.0 International licence, http://creativecommons.org/licenses/by/4.0/, as recorded in the arXiv licence metadata for this exact version and shown on the abstract page https://arxiv.org/abs/2510.12244v5. Full licence notice: \texttt{licenses/arxiv-2510.12244-CC-BY-4.0.txt}.\par\smallskip

\noindent\textit{Editorial scope.} Counted unit: the article's theorem loc:completeness\_of\_the\_composition\_of\_nested\_normals.statement (source0.tex lines 933--940). The article's complete original proof of it is delivered with it (source0.tex lines 960--996, one \texttt{proof} environment of 37 lines, which the article places away from the statement). No mathematics is written, repaired or re-derived: the statement and the proof are the article's own lines, in the article's own notation, and no retained line is substituted or re-flowed.  Retained uncounted as the source-local context the proof needs: lines 725--738; lines 946--951.  Nothing was omitted from any retained span, so the package carries no omission record.  No retained line contains a citation, so the wrapper carries no bibliography and no bibliography entry.  No retained line contains a figure environment, an image-inclusion command or drawing code, so no image is delivered and the package declares no asset dependency.  Editorial formatting only, in addition: an AMS \texttt{amsart} preamble that restates the article's own declarations and macros used by the retained lines, namely the environments \texttt{proposition}, \texttt{theorem}, \texttt{lemma} on separate counters, together with the standard packages those lines need.  The retained environments print in this document as Proposition 1, Theorem 1, Lemma 1 in order of appearance, whereas the article prints the same items with the numbering of its own full document.  The complete native source of the article as distributed in the arXiv e-print for this exact version is bundled unchanged as \texttt{sources/187026/source0.tex}.
\begin{proposition}[Relation of supporting subspaces to faces]
\label{loc:relation_of_supporting_subspaces_to_faces.statement}
Let $C \subseteq X$ be convex.
\begin{enumerate}
\item For any affine subspace $L \subseteq X$,
\begin{equation*}
L \text{ is a supporting subspace }  \iff L \cap C \text{ is a face.}
\end{equation*}
\item For any subset $F \subseteq C$,
\begin{equation*}
F \text{ is a face of }C  \iff \exists L \text{ a supporting subspace s.t.} F = L \cap C.
\end{equation*}
\end{enumerate}
\end{proposition}

\begin{lemma}[Affine slices preserve supporting subspaces]
\label{loc:affine_slices_preserve_supporting_subspaces.statement}
Let $C \subseteq \mathbb{R}^n$ be convex and $S \subseteq C$ non-empty.
Let $U \in \mathcal{F}(C;S)$, and $V$ any subspace such that
$S-S \subseteq V$. Then $U \cap V \in \mathcal{F}(C \cap (V+S);S)$.
\end{lemma}

\begin{theorem}[Completeness of the composition of nested normals]
\label{loc:completeness_of_the_composition_of_nested_normals.statement}
Let $C \subseteq \mathbb{R}^n$ be convex, $S \subseteq C$ non-empty,
and $V, U \in \mathcal{F}(C;S)$ such that $U \subseteq V$.
Then there is a sequence $V_0, \ldots, V_k \in \mathcal{F}(C;S)$ 
with $k= \dim(V)-\dim(U)$ such that $V_0=V$, $V_k=U$, and for each
$i \in \{ 1, \ldots,k\}$ there is a $v_i \in N_C^{V_{i-1}}(S) \setminus \{0\}$ such that $V_i = V_{i-1} \cap v_i^\perp$.
\end{theorem}

\begin{proof}[\hypertarget{loc:completeness_of_the_composition_of_nested_normals.proof}Proof of \Cref{loc:completeness_of_the_composition_of_nested_normals.statement}]

It suffices to show the following statement: given any $U,V \in \mathcal{F}(C;S)$ such that $V \supsetneq U$,
there exists a $v \in N_C^V(S) \setminus \{0\}$ such that $v \perp U$. 
If this statement is shown to hold, applying it recursively generates the sequence of supporting subspaces $V_1, \ldots, V_k$ in a finite number of steps. 
Note that the dimension of the supporting subspaces decreases by one at each iteration while always containing $U$, which implies that $U \subseteq V_k$ with $\dim(U)=\dim(V_k)$, and therefore $U=V_k.$

We now show the above statement.
Let $C' = C \cap (V + S)$, $\bar{C} = \proj_{U^\perp} C'$, and $\bar{U} =  \proj_{U^\perp} (U + S)$.
Note that $\bar{C}$ is a convex set and $\bar{U}$
a singleton $\{x\} \subseteq \bar{C}$.  We argue that $\bar{U}$ is a face of $\bar{C}$.
Arguing by contradiction, we hypothesize the $\bar{U}$ is not a face of $\bar{C}$, in which case there are $\bar{y}, \bar{z} \in \bar{C}\setminus \bar{U}$
such that $x = \lambda\bar{y} + (1-\lambda)\bar{z}$ for $\lambda \in (0, 1)$. This implies that there are elements $y, z \in C'$ such that $\bar{y}= \proj_{U^\perp}y$ 
and $\bar{z}= \proj_{U^\perp}z$, with $w =  \lambda y+(1-\lambda)z$. Note that the line segment $[y,z]$ has endpoints $y, z \in C' \setminus (U+S)$ with a point $w  \in ]y,z[$  that is also contained in $C' \cap (U+S)$. This implies that $C' \cap (U+S)$ is not a face of $C'$, which, by \Cref{loc:relation_of_supporting_subspaces_to_faces.statement}, means that
$U \notin \mathcal{F}(C';S)$. This is a contradiction because in fact, $U  \in \mathcal{F}(C';S)$, as we now show. We know that $U  \in \mathcal{F}(C;S)$, so by
\Cref{loc:affine_slices_preserve_supporting_subspaces.statement}, $U = U \cap V \in \mathcal{F}(C \cap (V+ S);S)$, i.e., $U \in \mathcal{F}(C';S)$.
Therefore, we indeed have a contradiction, from which we find that $\bar{U}$ is a face of $\bar{C}$.

Note that $\proj_{U^\perp}(V + S) \subseteq V + S$. This becomes
apparent from the following equality:
$V + S = (V + S) + U = \proj_{U^\perp} (V + S) + U,$
where the first equality holds because $U \subseteq V$ implies that $V = V + U$. 

Then since $U+S, C \cap (V +S ) \subseteq V + S$, we find that $\bar{U}, \bar{C} \subseteq (V+S) \cap U^\perp$.
Since $\bar{U}$ is a face of $\bar{C}$, 
there exists a supporting hyperplane $H$ of $\bar{C}$, where
the hyperplane is relative to the affine subspace $(V+S) \cap U^\perp$, such that
$\bar{U} \subseteq H$.
Indeed, if there is no such hyperplane $H$, $\bar{U}$ would then be a singleton
in the interior of $\bar{C}$ (where the interior is relative to $(V+S)\cap U^\perp$, i.e., with the
induced topology in $(V+S)\cap U^\perp$).

Take $\bar{v} \in V \cap U^\perp$ to be the unit
normal to $H$ that is also in $N_{\bar{C}}^{V \cap U^\perp}(x)$. 
Then note that $\bar{v} \in N_C^{V}(S) \cap U^\perp$, which is the vector
we had to find.
\end{proof}

\end{document}
